\documentclass[a4paper,12pt]{article} 
\usepackage{mathrsfs}
\usepackage[utf8]{inputenc}
\usepackage[spanish]{babel}
\usepackage{amsmath}
\usepackage{amsfonts}
\usepackage{amssymb} 
\usepackage{graphicx} 
\usepackage{hyperref} 
\usepackage{wrapfig}
\usepackage{enumitem}
\usepackage{fancyhdr}
\usepackage{float}
\usepackage{eurosym}
\usepackage{color}
\usepackage{circuitikz}
\usepackage{titling}
\usepackage{hyperref}
\usepackage{media9}
\usepackage{lipsum}
\usepackage{tocbibind}
\usepackage{listings}
\usepackage{tabularx}
\usepackage{tcolorbox}
\usepackage{bookmark}
\usepackage{media9}
\usepackage[table]{xcolor}
\definecolor{lightblue}{RGB}{228, 244, 253}
\usepackage{color}

\definecolor{dkgreen}{rgb}{0,0.6,0}
\definecolor{gray}{rgb}{0.5,0.5,0.5}
\definecolor{mauve}{rgb}{0.58,0,0.82}

\lstset{frame=tb,
  language=Python,
  inputencoding=utf8,
  extendedchars=true,
  aboveskip=3mm,
  belowskip=3mm,
  showstringspaces=false,
  columns=flexible,
  basicstyle={\small\ttfamily},
  numbers=none,
  numberstyle=\tiny\color{gray},
  keywordstyle=\color{blue},
  commentstyle=\color{dkgreen},
  stringstyle=\color{mauve},
  breaklines=true,
  breakatwhitespace=true,
  tabsize=3,
  literate=%
    {á}{{\'a}}1
    {é}{{\'e}}1
    {í}{{\'i}}1
    {ó}{{\'o}}1
    {ú}{{\'u}}1
    {ñ}{{\~n}}1
    {č}{{\v{c}}}1
}
\usepackage[left=3cm,right=3cm,top=3cm,bottom=4cm]{geometry}
\sloppy

\pagestyle{fancy}
\providecommand{\abs}[1]{\lvert#1\rvert}
\providecommand{\norm}[1]{\lVert#1\rVert}

%%% Para las cabeceras
\newcommand{\hsp}{\hspace{20pt}}
\newcommand{\HRule}{\rule{\linewidth}{0.5mm}}
\headheight=50pt
%%% 
\newcommand{\vacio}{\textcolor{white}{ .}}

%%% Para que las ecuaciones se numeren
%%% con el número de sección y el de ecuación.
\renewcommand{\theequation}{\thesection.\arabic{equation}}


% Color azul para algunos 
% textos de la portada
\definecolor{azulportada}{rgb}{0.16, 0.32, 0.75}

%%%% Azul para textos de headings
\definecolor{azulinterior}{rgb}{0.0, 0.2, 0.6}

%%%%%%%%%%%%%%%%%%%%%%%%%%%%%%%%
%%%%%% Datos del proyecto %%%%%%
%%%%%%%%%%%%%%%%%%%%%%%%%%%%%%%%
%%%TÍTULO
%%% Escribirlo en minúsculas, el programa
%%% lo pondrá en mayúsculas en la portada.

\title{Primeros pasos con ROS 2}

%%%% AUTORES
\author{Lydia Ruiz Martínez \and Pablo Tuñón Laguna \and Jorge Vančo Sampedro}

%%%%%%%%%%%%%%%%%%%%%
%%%%%%%%%%%%%%%%%%%%
\begin{document}

%%%%%%%%%%%%%%%%%%%%%%%%%%%%%%%
%%%%%%%%%%%%%%%%%%%%%%%%%%%%%%%
\begin{titlepage} %%%%% Aquí no hay que tocar nada.
	%%%% Las siguientes instrucciones generarán automáticamente
	%%%% la portada de tu proyecto.
	%%% Cambio de la estructura de esta página
\newgeometry{left=0.6cm,top=1.3cm,bottom=1.2cm}

\fbox{\parbox[c]{18.5cm}{
\begin{center}
\vspace{1.5cm}
{\fontfamily{ptm}\fontsize{24}{28.8}\selectfont{Universidad Pontificia Comillas}}\\
[2.5em]
{\fontfamily{ptm}\fontsize{24}{5}\selectfont{ETSI-ICAI}}\\
[3.5em]
{\fontfamily{ptm}\fontsize{28}{5}\selectfont{Laboratorio Robots autónomos}}\\
[3cm]
% {\fontfamily{ptm}\fontsize{24}{5}\selectfont{Visión por Ordenador I}}\\
% [2cm]

% Autor del trabajo de investigación
\textcolor{azulportada}
{\fontfamily{ptm}\fontsize{45}{5}\selectfont{\textsc{\thetitle}}}\\
[3cm]
% Título del trabajo
\textcolor{azulportada}
{\fontfamily{ptm}\fontsize{16}{5}\selectfont{\theauthor}}\\
%{\Huge\textbf{\thetitle}}\\
[1.2cm]
\includegraphics[width=10cm]{fonts/Logo ICAI.png}
\\[1.8cm]

{\fontfamily{ptm}\fontsize{16}{5}\selectfont{December 2024}}\\
[4cm]
\end{center}
}}
\end{titlepage}
 
 \restoregeometry
 %%%% Volvemos a la estructura de la página normal

%%%%%%%%%%%%%%%%%%%%%%%%%%%%%%

{%\Large

%%%Encabezamiento y pie de página
%%% También se genera automáticamente
%%% Mejor no tocarlo mucho.
\renewcommand{\headrulewidth}{0.5pt}
\fancyhead[R]{
\textcolor{azulportada}{\fontfamily{ptm}\fontsize{10}{3}\selectfont{Proyecto Final de Visión por Ordenador I}}\\
{\fontfamily{ptm}\fontsize{10}{3}\selectfont{\theauthor}}}
\fancyhead[L]{
  \textcolor{azulinterior}{\fontfamily{ptm}\fontsize{12}{4}\selectfont{\textbf{\thetitle}}}\\
}


\pagestyle{fancy}
\renewcommand{\footrulewidth}{0.5pt}
\fancyfoot[L]{\footnotesize Universidad Pontificia Comillas (ICAI) --- curso 2024-2025}
\fancyfoot[C]{\vacio}
\fancyfoot[R]{\footnotesize Página \thepage}

%%%%%%%%%%%%%%%%%%%%

% \newpage

% \renewcommand{\contentsname}{Índice de contenidos}
% \addtocontents{toc}{\protect\setcounter{tocdepth}{-1}} % Quita el índice de la tabla de contenidos
% \tableofcontents
% \addtocontents{toc}{\protect\setcounter{tocdepth}{3}}

\newpage

\textbf{¿Cuál es la versión y el nombre de la distribución de ROS instalada en el contenedor que se le ha proprocionado? 
¿Cuándo finaliza su soporte? ¿Por qué?}


\vspace{1cm}

La versión de ROS instalada es la 2a, la distribuciones usadas son Humble y Jazzy. El soporte de Humble finaliza en mayo de 2027 
y el de Jazzy en mayo de 2029 (en el día internacional de las tortugas). Ambas distribuciones surgieron en años pares (2022 y 2024 respectivamente)
por lo que poseen el soporte extendido de 5 años.

\vspace{0.5cm}

\begin{lstlisting}
  printenv ROS_VERSION # Extrae la versión de ROS
  printenv ROS_DISTRO # Extrae la distribución de ROS
\end{lstlisting}

\vspace{1cm}


\textbf{¿Qué otro comando se puede ejecutar con ros2 node?}


\vspace{1cm}

El primer comando inicializa un nodo de ROS2 que contiene información de una tortuga. El segundo comando 
muestra los nodos que se encuentran actualmente en ejecución y el tercer comando muestra información sobre el nodo de la tortuga.

\vspace{0.5cm}

\begin{lstlisting}
  ros2 run turtlesim turtlesim_node # Inicializa un nodo de ROS2
  ros2 node list # Muestra los nodos en ejecución
  ros2 node info /turtlesim # Muestra información sobre el nodo de la tortuga
\end{lstlisting}

\vspace{1cm}


\textbf{¿Cómo se llama el tema (topic) donde /teleop\_turtle publica los comandos de movimiento?}


\vspace{1cm}

El tema se llama /turtle1/cmd\_vel, obtenido a través de la herramienta rqt\_graph.

\vspace{0.5cm}

\begin{lstlisting}
  rqt_graph # Muestra un gráfico de nodos y temas en ROS2
\end{lstlisting}

\vspace{0.5cm}

\begin{figure}[H]
    \centering
    \includegraphics[width=0.8\textwidth]{fonts/rqt_graph_1.png}
    \caption{Gráfico de nodos y temas en ROS2 usando rqt\_graph}
    \label{fig:rqt_graph}
\end{figure}

\vspace{1cm}


\textbf{¿Qué recibe /michelangelo de /teleop\_turtle?}


\vspace{1cm}

No recibe nada ya que no todavía no se ha publicado nada, cuando se va a la terminal del nodo teleoperador y se pulsa una tecla, se muestra el 
mensaje que se envía a través del topic. El mensaje es de tipo geometry\_msgs/msg/Twist y contiene la velocidad lineal y angular de la tortuga.

\vspace{0.5cm}

\begin{lstlisting}
  ros2 topic echo /turtle1/cmd_vel # Muestra los mensajes publicados en un topic
  ros2 topic info /turtle1/cmd_vel # Muestra información sobre un topic
\end{lstlisting}

\vspace{1cm}


\textbf{¿Qué tipos de datos se intercambian entre el nodo de teleoperación y TurtleSim?}


\vspace{1cm}

Se intercambian 2 topics con distintos mensajes y una acción:

\vspace{0.5cm}

\begin{itemize}
    \item /turtle1/cmd\_vel: Mensaje de tipo geometry\_msgs/msg/Twist que contiene la velocidad lineal y 
    angular de la tortuga como vectores de 3 elementos float64 (x, y, z).
    \item /turtle1/rotate\_absolute/\_action/status: Mensaje de tipo action\_msgs/msg/GoalStatusArray que 
    contiene el estado de la acción como un array de enteros.
    \item /turtle1/rotate\_absolute/\_action/feedback: acción de tipo Feedback.
\end{itemize}

\vspace{0.5cm}

\begin{lstlisting}
  ros2 topic info /turtle1/cmd_vel # Muestra información sobre un topic
  ros2 interface show geometry_msgs/msg/Twist # Muestra información sobre un mensaje
  ros2 topic info /turtle1/rotate_absolute/_action/status # Muestra información sobre un topic
  ros2 interface show action_msgs/msg/GoalStatusArray # Muestra información sobre un mensaje
\end{lstlisting}

\vspace{1cm}


\textbf{¿Qué campos hay que completar en package.xml? ¿Qué es exactamente una licencia en este contexto? 
¿Puede citar algunas licencias de software?}


\vspace{1cm}

Hay que completar:

\vspace{0.5cm}

\begin{itemize}
  \item description: Descripción del paquete.
  \item maintainer: Nombre y correo electrónico del mantenedor del paquete.
  \item license: Licencia del paquete.
  \item exe\_depends: Dependencias del paquete, rclpy y std\_msgs. 
\end{itemize}

\vspace{0.5cm}

En este contexto, una licencia es un contrato que establece los términos y condiciones bajo los cuales se puede usar, 
modificar y distribuir un software. En concreto la Apache 2.0 es una licencia de software libre y de código abierto 
que permite a los usuarios utilizar, modificar y distribuir el software siempre y cuando se de credito al autor 
original y se incluya la licencia en cualquier copia o modificación del software. Otras licencias son la GPL, MIT, BSD.

\vspace{1cm}


\textbf{¿Qué hay que añadir a setup.py?}


\vspace{1cm}

Hay que añadir:

\vspace{0.5cm}

\begin{itemize}
    \item maintainer: Nombre y correo electrónico del mantenedor del paquete.
    \item maintainer\_email: Correo electrónico del mantenedor del paquete.
    \item description: Descripción del paquete.
    \item license: Licencia del paquete.
    \item entry\_points: Punto de entrada del paquete (desde donde se ejecuta) tanto para el publisher como para el subscriber.
\end{itemize}

\vspace{1cm}


\textbf{¿Cuáles son los pasos para construir y ejecutar un nodo?}


\vspace{1cm}

Los pasos son:

\vspace{0.5cm}

\begin{itemize}
    \item Crear un workspace de ROS2.
    \item Crear un paquete.
    \item Diseñar los nodos y mensajes necesarios.
    \item Modificar el package.xml para añadir el mantenedor,  la descripción, la licencia y las dependencias.
    \item Modificar el setup.py para añadir el mantenedor, la descripción, la licencia y el punto de entrada.
    \item Compilar el paquete.
    \item Ejecutar el nodo.
\end{itemize}

\vspace{0.5cm}

\begin{lstlisting}
  mkdir -p ~/ros2_ws/src # Crea un workspace de ROS2
  cd ~/ros2_ws/src
  ros2 pkg create --build-type ament_python <package_name> # Crea un paquete
  # Programar los nodos y mensajes necesarios
  # Modificar package.xml y setup.py
  colcon build # Compila el paquete
  source ~/ros2_ws/install/setup.bash # Añade el paquete al workspace
  ros2 run <package_name> <node_name> # Ejecuta el nodo
\end{lstlisting}

\vspace{1cm}


\textbf{En sus propias palabras, explique brevemente cómo se comunica el nodo parlante
(talker) con el oyente (listener)}


\vspace{1cm}

El nodo publisher, publica un mensaje en un tópic (un canal de comunicación), el nodo subscriber está suscrito a dicho topic
de tal manera que cuando se publica un mensaje en dicho canal, el nodo subscriber recibe el mensaje y lo procesa.

\vspace{1cm}


\textbf{Utilice las herramientas de terminal de ROS 2 y rqt para descubrir los nodos que
operan el TurtleBot3, así como los publicadores y suscriptores que tienen. Explore también la estructura
y el contenido de los mensajes e intente explicar con sus palabras para qué sirve cada nodo a partir de
esa información. Ignore todos los temas (topics) que empiecen con el prefijo rcl\_}


\vspace{1cm}

Bla, bla bla...

\vspace{1cm}


\textbf{¿Qué criterio de signos utiliza el TurtleBot3 para la velocidad angular $\omega$ por defecto? ¿Coincide con el habitual?}


\vspace{1cm}

Bla, bla bla...

\vspace{1cm}


\textbf{Las tortugas bebé corren hacia el mar}

\end{document}



